\subsection{The first nonabelian inert Frobenius and its exact Frattini orbit}

\begin{statement}{E.59}{Proved}{A uniform cubic-residue formula and the rank-five Frobenius orbit at seventeen}
Let $K=\mathbb Q(\zeta_3)$ and
$P=G_{K,\{3\}}^{(3)}$. Use the actual free basis of
E.55, written $x,y$, with
\[
 \chi(x)=4,\quad\kappa_3(3)(x)=0,\qquad
 \chi(y)=1,\quad\kappa_3(3)(y)=1,\qquad
 cxc=x,\quad cyc=y^{-1}.
 \tag{E59.1}
\]
Here $x$ is the specified arithmetic Frobenius at the
prime over two, and $c$ is the specified complex
conjugation in $P\rtimes C_2$.
Put $\Phi=P^3[P,P]$, $\Phi_2=\Phi^3[\Phi,\Phi]$,
and let $D_n(P)$ be the three-Zassenhaus filtration.
The Frattini field is
\[
 K_1=K(\zeta_9,\sqrt[3]{3}),\qquad
 \operatorname{Gal}(K_1/K)=(\mathbb Z/3)^2,\qquad
 \dim_{\mathbb F_3}\Phi/\Phi_2=10.
 \tag{E59.2}
\]
Every inert rational prime $\ell\equiv8\pmod9$
has arithmetic Frobenius $\sigma_\ell\in D_3(P)$.
The following explicit formula determines its
whole image in $D_3(P)/D_4(P)$.

Set $z=\zeta_9$, $\omega=z^3$, $r=\sqrt[3]{3}$,
and define actual elements of $K_1$ by
\[
 \begin{gathered}
 u=-z^2(1-z)^2,\quad
 u_1=-z^8(1-z^4)^2,\quad
 u_2=-z^{14}(1-z^7)^2,\\
 h=\frac{(r-u)(\omega^2r-u_1)}
          {(\omega^2r-u)(r-u_1)}.
 \end{gathered}
 \tag{E59.3}
\]
The denominators have no prime divisor away
from three. In $\mathbb F_{\ell^2}$ choose a
primitive ninth root $\bar z$, put
$\bar\omega=\bar z^3$, and take the unique
cube root $\bar r$ of $3$ in $\mathbb F_\ell$.
Define $\alpha_\ell,\beta_\ell\in\{0,1,2\}$ by
\[
 \alpha_\ell=\frac{\ell^2-1}{9}\pmod3,\qquad
 \bar h^{\,(\ell^2-1)/3}=\bar\omega^{\,\beta_\ell}.
 \tag{E59.4}
\]
Then the actual comparison is
\[
 \boxed{\quad
 \sigma_\ell
   =x^{3\alpha_\ell}[y,[x,y]]^{\beta_\ell}
                  \quad\text{in }P/D_4(P).
 \quad}
 \tag{E59.5}
\]
Changing the prolongation, or either of the
compatible residue-field roots in (E59.4), gives
the same two coefficients.
For example,
\[
 \sigma_{17}=x^6[y,[x,y]]^2\pmod{D_4(P)},\qquad
 \sigma_{17}\in\Phi\setminus\Phi_2,\qquad
 \sigma_{17}\in D_3(P)\setminus D_4(P).
 \tag{E59.6}
\]

Here is an integral-source comparison which
also determines the exact span of its conjugates.
Write
\[
 R=\mathbb F_3[X,Y]/(X^3-1,Y^3-1)
   =\mathbb F_3[s,t]/(s^3,t^3),\qquad
 s=X-1,\quad t=Y-1,\quad I=(s,t).
 \tag{E59.7}
\]
Continuous Fox derivatives give an actual
$R$-module isomorphism
\[
 \Phi/\Phi_2
       \xrightarrow{\ \sim\ }
 M=\ker\bigl(R^2\xrightarrow{(s,t)}R\bigr).
 \tag{E59.8}
\]
The involution induced by $c$ is
\[
 (A,B)\longmapsto
  \bigl(A(X,Y^{-1}),-Y^{-1}B(X,Y^{-1})\bigr).
 \tag{E59.9}
\]
Its fixed subspace has the exact basis
\[
 a=(s^2,0),\qquad
 b=(-t^2,st-st^2),\qquad sb,\qquad s^2b.
 \tag{E59.10}
\]
For
$w=Aa+(B+Cs+Ds^2)b$, its entire cyclic
$R$-module has dimension
\[
 \begin{array}{c|c|c}
 \text{coordinates}&\dim_{\mathbb F_3}Rw&
                       \text{one basis}\\ \hline
 B\ne0&5&w,sw,tw,s^2w,stw\\
 B=0,\ A\ne0&3&w,tw,t^2w\\
 A=B=0,\ C\ne0&3&w,sw,tw\\
 A=B=C=0,\ D\ne0&1&w\\
 A=B=C=D=0&0&\varnothing .
 \end{array}
 \tag{E59.11}
\]
A prolongation with $c\sigma_\ell c=\sigma_\ell$
exists, and its coordinates in (E59.10) satisfy
$A=\alpha_\ell$, $B=\beta_\ell$.
Thus the nine conjugate Frobenius classes have
rank five when $\beta_\ell\ne0$, and rank three
when $\beta_\ell=0$, $\alpha_\ell\ne0$.
In particular, their rank at seventeen is
\emph{exactly five}. Each individual class at
seventeen is nonzero, but the nine classes
cannot be part of a free generating set of
$\Phi$. This last conclusion follows from
the displayed map on its Frattini quotient.

For completeness, here is a separate full
arithmetic certificate of the rank at seventeen.
In the ten-dimensional Kummer basis
$e_1,\ldots,e_8,\pi_3,z$ specified by the
retained exact coefficient file, the nine
Frobenius columns are
\[
 \begin{pmatrix}
 1&1&1&1&1&1&1&1&1\\
 0&0&0&0&0&0&0&0&0\\
 1&1&1&2&2&2&0&0&0\\
 1&1&1&0&0&0&2&2&2\\
 1&1&1&1&2&0&0&1&2\\
 2&0&1&1&0&2&1&1&1\\
 0&0&0&1&0&2&0&2&1\\
 1&0&2&0&0&0&2&1&0\\
 2&0&0&2&2&1&0&2&0\\
 2&2&2&2&2&2&2&2&2
 \end{pmatrix}.
 \tag{E59.12}
\]
Its rank is five; its kernel has dimension four
and contains the all-one vector.
The residue procedure is evaluated at every
one of the fifty primes $\ell\le2000$ in this
congruence class. The first coefficients include
\[
 \begin{array}{c|rrrrrrrrrr}
 \ell&17&53&71&89&107&179&197&233&251&269\\ \hline
 \alpha_\ell&2&0&2&1&0&2&1&2&1&0\\
 \beta_\ell &2&1&1&0&2&1&2&2&2&0 .
 \end{array}
 \tag{E59.13}
\]
The three primes in this range with both
coefficients zero are $269,809,1187$.
Their full ten-unit residue matrices, also
retained, have respective ranks $3,1,3$.
These ranks use the entire Frattini quotient,
including the coordinates beyond (E59.5).
There is also an exact distribution theorem.
For ranks $0,1,3,5$, respectively, the relative
Dirichlet densities among rational primes
$\ell\equiv8\pmod9$ are
\[
       \frac1{81},\qquad\frac2{81},\qquad
       \frac8{27},\qquad\frac23.
 \tag{E59.13a}
\]
Their absolute Dirichlet densities are
$1/486,1/243,4/81,1/9$.
They follow from the actual finite Galois
quotient below.

The maps in this statement are the actual
maps
$P\to P/\Phi$, $\Phi\to\Phi/\Phi_2$ and
$P\to P/D_4(P)$, together with their Kummer
duals. The correction words in the arithmetic
inertia presentation of E.39 are retained.
The calculation does not replace an actual
local Frobenius by the image of a chosen
section of an arithmetic quotient.
The manuscript anchor is \S51.16, folios
733v--736v, scans~172--178, formulas~(371)--(379),
and Part~E(c),(d).
\end{statement}

\begin{proof}
\emph{The arithmetic field and the two filtrations.}
E.34 and the actual basis construction in E.55
give (E59.1). The two independent mod-three
characters are the cyclotomic character on
$\zeta_9$ and the Kummer character of $3$.
Their fixed field is exactly $K_1$.
The subgroup $\Phi$ has index nine in the
free pro-three group of rank two, so its
Schreier rank is $1+9(2-1)=10$.

For $\ell\equiv8\pmod9$, cubing is an
automorphism of $\mathbb Q_\ell^\times$ on
its unit subgroup: both its residue-unit
group and its principal-unit group have
order prime to three at every finite
quotient. All compatible cube-power roots
of $3$ can therefore be taken in
$\mathbb Q_\ell$. The local Frobenius has
$\chi=\ell^2$ and $\kappa_3(3)=0$.
It is trivial on both generators of $K_1/K$,
so it lies in $\Phi=D_2(P)$.

There is also a precise normalization by
complex conjugation. The image of the local
decomposition group in $P\rtimes C_2$ is
procyclic, since the extension is unramified
at $\ell$. Its cyclotomic image identifies
it with $C_2\times\mathbb Z_3$. Its unique
involution is conjugate to $c$ by an element
of $P$, by the coprime complement theorem.
After making this actual conjugation, its
Frobenius commutes with $c$. Both involutions
have zero Kummer character on $3$; the
conjugating element has zero such character
because the conjugation formula gives twice
that character. Hence the normalization also
preserves $\kappa_3(3)(\sigma_\ell)=0$.

The quotient $D_2/D_3$ is generated by
$[x,y]$ and $c$ acts on it by minus one.
Consequently the $c$-fixed Frobenius lies
in $D_3$. Collection, or the truncated
Magnus expansion over $\mathbb F_3$, gives
the four independent generators
\[
 x^3,\quad y^3,\quad[x,[x,y]],\quad[y,[x,y]]
                    \quad\text{of }D_3/D_4.
 \tag{E59.14}
\]
Their signs under $c$ are $+,-,-,+$.
Thus only the first and the last can occur.
On $x^3$ the cyclotomic character is $64$,
and $(64-1)/9\equiv1\pmod3$. This proves
the formula for $\alpha_\ell$.
At seventeen $\chi(\sigma_{17})=289$ is
one modulo nine but not modulo twenty-seven.
Both $\Phi_2$ and $D_4$ have cyclotomic
image contained in $1+27\mathbb Z_3$.
This proves the two strict assertions
in (E59.6).

\emph{The Fox map, including the fixed lattice.}
The degree-nine covering of the two-circle
graph with deck group $P/\Phi$ has chain
complex $R^2\to R$ with boundary $(s,t)$.
Its first homology is the Frattini quotient
of its free pro-three fundamental group.
This gives (E59.8), and also proves its
compatibility with conjugation by $x,y$.
The Fox derivatives of $x$ and $y^{-1}$
are respectively $1$ and $-Y^{-1}$.
They give (E59.9) without a change of
coefficient convention.

Write a polynomial in $t$ to degree two.
The first coordinate of a fixed vector has
the form $A_0(s)+A_2(s)t^2$, while the
second has the form $B_1(s)(t-t^2)$.
The boundary equation is
$sA_0=0$ and $sA_2+B_1=0$.
It follows that $A_0=As^2$ and
$A_2=-(B+Cs+Ds^2)$.
This proves the complete basis (E59.10).
The leading Fox coordinates of $x^3$ and
$[y,[x,y]]$ are respectively
$(s^2,0)$ and $(-t^2,st)$.
Thus the first two fixed coordinates are
exactly the two coefficients in (E59.5).

Here are all the cyclic-module calculations
in (E59.11). For $B\ne0$, the five displayed
vectors have independent lowest homogeneous
terms. The ideal $I^3$ and the polynomial
$As^2+Bt^2$ annihilate $w$; their quotient
in $R$ has dimension five. For $B=0,A\ne0$,
$s+(C/A)t^2$ annihilates $w$, and
$w,tw,t^2w$ are independent. For
$A=B=0,C\ne0$, $I^2$ annihilates $w$,
and $w,sw,tw$ are independent.
The remaining two cases follow by
multiplying $s^2b$ by $s$ and $t$.
These identities prove the complete table,
including all degenerate cases.
The retained finite program verifies
all eighty-one fixed vectors as a second
certificate.

There is a universal relation as well.
For a cycle $(A,B)$, reducing
$sA+tB=0$ modulo $I^2$ shows
$A,B\in I$. The group norm of
$(\mathbb Z/3)^2$ is $s^2t^2$; it
annihilates $I$. Thus the sum of all
nine conjugates of every element of
$\Phi/\Phi_2$ is zero.
The rank-five assertion at seventeen
is stronger than this one relation.

\emph{The finite arithmetic quotient and the density count.}
Let $G_0=P/\Phi_2$. The exact Fox map embeds
it as the following entire set of pairs:
\[
 \begin{split}
 G_0&=\{(q,D):q\in\langle X,Y\rangle,\quad
          D\in R^2,\quad(s,t)D=q-1\},\\
 (q,D)(q',D')&=(qq',D+qD').
 \end{split}
 \tag{E59.13b}
\]
Indeed, the Fox identity gives the boundary
condition. For each fixed $q$, a chosen word
over $q$ supplies one derivative, and
(E59.8) supplies every difference in $M$.
The kernel is exactly $\Phi_2$.
Thus $|G_0|=3^{12}$.
The characteristic subgroup $\Phi_2$ is
stable under $c$, so
$G=G_0\rtimes\langle c\rangle$ is an
actual finite Galois quotient over
$\mathbb Q$, unramified outside three.

The cyclotomic condition
$\chi_9=-1$ is precisely the subset
$g=(Y^b,D,c)$, with $b=0,1,2$ and
$(s,t)D=Y^b-1$. Its size is $3^{11}$.
Write $D=D_b+m$, where $D_b$ is the
Fox derivative of $y^b$ and $m\in M$.
The square $g^2$ is in $M$ and is
\[
 T_b(D)=D+Y^b c(D),\qquad
 T_b(D_b)=0,\qquad
 Y^bT_b(D)=(1+c)Y^b m .
 \tag{E59.13c}
\]
The last equality uses $Y^3=1$ and
the actual semilinear involution
(E59.9). Multiplication by $Y^b$
is an invertible $R$-module map.
Since two is invertible in $\mathbb F_3$,
$(1+c)Y^b:M\to M^c$ is surjective,
with kernel of dimension six.
Every one of the eighty-one fixed
vectors therefore has $3^7$ preimages
in the complete cyclotomic subset.

The five cases in (E59.11) give exactly
$1,2,24,54$ fixed vectors of ranks
$0,1,3,5$. Their corresponding finite
Galois subsets have sizes
$2187,4374,52488,118098$.
The rank is invariant under all
conjugations in $G$: conjugation by
$P$ multiplies the cycle by an
invertible deck element, and $c$
acts semilinearly. For a rational
prime in this cyclotomic class,
$g^2$ is its actual arithmetic
Frobenius over $K$.
Chebotarev's theorem for the resulting
conjugacy-stable subsets
\cite[VIII, Theorem 7.4]{E_MilneCFT}
therefore gives (E59.13a), upon
division by $|G|=2\cdot3^{12}$.
The retained program checks the
ten-generator Schreier basis, all
three maps in (E59.13c), and these
exact finite counts.

\emph{An explicit character dual to the triple bracket.}
Let $X$ act on $K_1$ by $z\mapsto z^4$,
$r\mapsto r$, and let $Y$ act by
$z\mapsto z$, $r\mapsto\omega r$.
These are exactly the restrictions of
$x,y$. Direct calculation in
$\mathbb Q[z]/(z^6+z^3+1)$ gives
\[
 u u_1u_2=3,\qquad
 \gamma=u_1u_2^2,\qquad
 q=\frac{-2z^5+2z^4-z^2+z}{3},
 \qquad
 \frac{3-u^3}{3-u_1^3}\,\gamma=9q^3.
 \tag{E59.15}
\]
The resultant of $3-u^3$ and $z^6+z^3+1$
is $3^{10}$.
Consequently each factor $r-u_i$ in
(E59.3), and each of its $Y$-conjugates,
is a unit away from three.
In particular $h$ defines a Kummer
character of $\Phi$ unramified outside
three.

Set $\phi=(r-u)/(r-u_1)$, so that
$h=\phi/(Y^{-1}\phi)$.
The exact norm identity (E59.15), with
$r^6=9$, gives the literal equality
\[
 \frac{Y^{-1}h}{h}
       =\gamma
          \left(\frac{Y^{-1}\phi}{r^2q}\right)^3.
 \tag{E59.16}
\]
Also $X\gamma/\gamma=(u/r)^3$ and
$Y\gamma=\gamma$. These are actual
element identities. In the Kummer group
they imply that $h$ is killed by
$(X-1)^2,(X-1)(Y-1),(Y-1)^2$:
the first follows also from $N_X\phi=1$,
and the last from $N_Yh=1$.
Furthermore $c(u_i)=\omega u_i$, so
$c(\phi)=Y^{-1}\phi$ and $c(h)=h^{-1}$.

Here is the exact filtration comparison
behind the character calculation.
The Fox module has dimension ten and
$I^2M$ has dimension five, by multiplying
the explicit cycles
$(-t,s),(s^2,0),(0,t^2)$.
We have $\Phi_2\subset D_4$ and
$I^2M\subset D_4/\Phi_2$.
The five independent collected classes
$[x,y],x^3,y^3,[x,[x,y]],[y,[x,y]]$
give $\dim\Phi/D_4=5$.
Hence
\[
                 D_4/\Phi_2=I^2M.
 \tag{E59.17}
\]
The Kummer character of $h$ therefore
factors through $\Phi/D_4$.
The sign $c(h)=h^{-1}$ combines with
$c(\omega)=\omega^{-1}$ to make that
character fixed under conjugation by $c$.

The element $\gamma$ is not a cube in $K_1$.
Indeed put $F=\mathbb Q(z)$. A cube root
in $F(r)$ would be an eigenvector for
$Y$, hence would have the form $r^k v$
with $v\in F$ and $k\in\{0,1,2\}$.
It would follow that $\gamma=3^k v^3$.
Then $X\gamma/\gamma=u^3/3$ would be
a cube in $F$, forcing $3$ to be a
cube there. But every intermediate
field of the abelian extension
$F/\mathbb Q$ is Galois over
$\mathbb Q$, whereas the real pure
cubic $\mathbb Q(\sqrt[3]{3})$
omits the two nonreal conjugates.
Thus the extension $K_1(\delta)$,
$\delta^3=\gamma$, has lifts
\[
 X(\delta)=(u/r)\delta,\qquad
 Y(\delta)=\delta .
 \tag{E59.18}
\]
Their commutator $[X,Y]$ sends
$\delta$ to $\omega\delta$; their
third powers are the identity.
Thus the pairing of $[x,y]$ with
$\gamma$ is exactly one, with this
commutator convention. For any Kummer
character $\kappa_h$ on $\Phi$,
conjugation gives
\[
 \kappa_h([y,g])
       =\kappa_{(Y^{-1}h)/h}(g)
             \quad(g\in\Phi).
 \tag{E59.19}
\]
Applying (E59.16) proves
$\kappa_h([y,[x,y]])=1$.

The other required value is
$\kappa_h(x^3)=0$.
At the chosen prime over two, $r$ reduces
to one, and the residue field in $K_1$
is $\mathbb F_{64}=\mathbb F_2[z]/(z^6+z^3+1)$.
Substitution in (E59.3) gives
$h=1+z^3=\omega^2$ there.
Its cubic residue is $h^{21}=1$.
By comparison, the character of $z$
on $x^3$ is one, and it vanishes on
every commutator. Thus $z,h$ are the
exact dual characters to
$x^3,[y,[x,y]]$ on the fixed part
of $D_3/D_4$.
The ordinary cubic-residue formula for
arithmetic Frobenius now proves
(E59.4)--(E59.5).
Conjugation by $P$ is trivial on
$D_3/D_4$, proving independence of
prolongation and of the associated roots.

At seventeen one can take
$\bar r=7$, $\bar z=2+8\bar\omega$,
where $\bar\omega^2+\bar\omega+1=0$.
Then $\bar u=6\bar\omega$,
$\bar u_1=12\bar\omega$,
$\bar h=15+9\bar\omega$, and
$\bar h^{96}=\bar\omega^2$.
Together with $(17^2-1)/9\equiv2\pmod3$
this proves the explicit word in (E59.6).

\emph{The complete residue certificates.}
The retained PARI source reconstructs
$K_1$ from $z^6+z^3+1$ and $r^3=3$,
with the actual primitive element
$\theta=r-z$. Its absolute polynomial is
\[
 \theta^{18}-21\theta^{15}+105\theta^{12}
 +326\theta^9+3633\theta^6+3660\theta^3+2197.
 \tag{E59.20}
\]
The unconditional certification procedure
of \cite{E_PARICertifiedUnits} returns one.
It certifies class number one, discriminant
$-3^{39}$, root group $\mu_{18}$, the
eight complete fundamental units, and
the unique prime above three.
The coefficient file supplies the literal
polynomials for these units, a generator
$\pi_3$ of that prime, $z,r$, the integral
basis, and every prime ideal used in
(E59.12). Thus the ten displayed
Kummer coordinates are specified actual
elements. Their classes form a basis
because valuations outside three are
divisible by three and the class group
is trivial. The subgroup $\Phi$ is the
maximal pro-three group unramified
outside three over $K_1$: taking a
normal closure of a finite three-extension
over the finite three-extension $K_1/K$
preserves this property.

At each of the nine degree-two primes
over seventeen, reduction of these ten
elements and raising to the ninety-sixth
power gives (E59.12), relative to $z^3$.
An independent finite-field implementation
matches the prime-ideal Hermite lattices
and verifies all ninety entries. It also
checks thirty residue entries at two,
the exact field actions, the low
Kummer subspace and the Heisenberg sign.
It independently evaluates (E59.4) for
all fifty relevant primes through two
thousand and for all nine prolongations,
giving four hundred fifty comparisons.
For $269,809,1187$ it evaluates all
ten Kummer characters, giving another
two hundred seventy residue values and
the exact ranks $3,1,3$.
These finite checks accompany the
separate finite-Galois proof of the
density theorem. The calculation
evaluates the displayed quotient maps;
it does not determine the full
Frobenius word in $P$.
\end{proof}
